\chapter{函数值与闭包}
\label{ch:15}

到目前为止，函数在你眼里还只是「一段有名字的代码」。可有时候，你并不关心那段代码叫什么，你只想把它本身当参数传出去——排序时告诉程序「怎么比」，遍历时告诉它「对每个元素做什么」。

这一章讲的就是这件事：函数也是一种值。而且它可以记住定义它的地方的一些变量——那就是闭包。

\section{问题：逻辑一样，只有「条件」不同}

写一个「从列表里筛出满足条件的元素」的函数。筛偶数是一个条件，筛大于 10 的是另一个条件：

\par
\medskip
\begin{lstlisting}[style=magel]
List<int64> filter(List<int64> xs, bool fn(int64) keep)
{
    List<int64> out;
    for (int64 x in xs)
    {
        if (keep(x))
        {
            out.push(x);
        }
    }
    return out;
}

int64 main()
{
    List<int64> xs;
    xs.push(3);
    xs.push(8);
    xs.push(15);
    xs.push(42);
    List<int64> evens = filter(xs, fn(int64 x)
    {
        x % 2 == 0
    });
    List<int64> big = filter(xs, fn(int64 x)
    {
        x > 10
    });
    println("{} {}", evens[0], evens[1]);
    println("{}", big.len());
    return 0;
}
\end{lstlisting}
\par
\medskip

\par
\medskip
\begin{lstlisting}[style=out]
8 42
2
\end{lstlisting}
\par
\medskip

\texttt{filter} 只管「遍历、判断、收集」，判断的标准由第二参数 \texttt{keep} 带进来。\texttt{keep} 的类型是 \texttt{bool fn(\allowbreak{}int64)\allowbreak{}}——「吃一个 \texttt{int64}、吐一个 \texttt{bool} 的函数」。传进去的两个 \texttt{fn(\allowbreak{}…)\allowbreak{}} 就是两个不同的条件。

一个 \texttt{filter}，无数种条件。 这就是「把函数当值」的价值。

\section{函数类型：\texttt{R fn(\allowbreak{}A,\allowbreak{} B)\allowbreak{}}}

先认这个写法：

\par
\medskip
\begin{lstlisting}[style=magel]
int64 fn(int64, int64)  // 吃两个 int64，返回 int64
bool fn(int64)  // 吃一个 int64，返回 bool
String fn()  // 不吃参数，返回 String
\end{lstlisting}
\par
\medskip

读法从左到右：返回类型 \texttt{fn(\allowbreak{}} 参数类型 \texttt{)\allowbreak{}}。函数类型也是类型，能写在变量、字段、参数、数组上——凡是你写 \texttt{int64}、\texttt{String} 的地方，都能写 \texttt{R fn(\allowbreak{}…)\allowbreak{}}。

\section{函数可以当值用}

把一个函数名赋给函数类型的变量，之后就能像调用原函数一样调它：

\par
\medskip
\begin{lstlisting}[style=magel]
int64 add(int64 a, int64 b)
{
    return a + b;
}
int64 sub(int64 a, int64 b)
{
    return a - b;
}

int64 apply(int64 fn(int64, int64) f, int64 a, int64 b)
{
    return f(a, b);
}

int64 main()
{
    int64 fn(int64, int64) op = add;
    println("{}", op(3, 4));
    op = sub;
    println("{}", op(3, 4));
    println("{}", apply(add, 20, 22));
    return 0;
}
\end{lstlisting}
\par
\medskip

\par
\medskip
\begin{lstlisting}[style=out]
7
-1
42
\end{lstlisting}
\par
\medskip

\texttt{op} 这个变量里装的是一个函数：

\begin{itemize}

\item \texttt{op =\allowbreak{} add;} 把 \texttt{add} 装进去；

\item \texttt{op =\allowbreak{} sub;} 换成 \texttt{sub}——同一个变量，指哪打哪；

\item \texttt{apply(\allowbreak{}add,\allowbreak{} 20,\allowbreak{} 22)\allowbreak{}} 把函数当普通实参传进去。

\end{itemize}

函数值能放的地方和你想象的一样多——字段、数组都行：

\par
\medskip
\begin{lstlisting}[style=magel]
struct Calculator
{
    int64 fn(int64, int64) op;
};

int64 add(int64 a, int64 b)
{
    return a + b;
}
int64 mul(int64 a, int64 b)
{
    return a * b;
}

int64 main()
{
    Calculator c;
    c.op = mul;
    println("{}", c.op(3, 4));

    int64 fn(int64, int64) table[2];
    table[0] = add;
    table[1] = mul;
    println("{} {}", table[0](6, 7), table[1](6, 7));
    return 0;
}
\end{lstlisting}
\par
\medskip

\par
\medskip
\begin{lstlisting}[style=out]
12
13 42
\end{lstlisting}
\par
\medskip

\texttt{table} 是一张「操作表」，按下标取函数来调——这是实现「操作码分派」最朴素的写法。

\section{\texttt{fn}：现场写一个小函数}

\texttt{add} / \texttt{sub} 都是提前定义好、有名字的函数。可很多时候，你只想当场写一小段逻辑，懒得给它起名。这就是 \texttt{fn}：

\par
\medskip
\begin{lstlisting}[style=magel]
int64 main()
{
    int64 fn(int64) square = int64 fn(int64 x)
    {
        x * x
    };
    println("{}", square(6));

    let cube = fn(int64 x)
    {
        x * x * x
    };
    println("{}", cube(3));
    return 0;
}
\end{lstlisting}
\par
\medskip

\par
\medskip
\begin{lstlisting}[style=out]
36
27
\end{lstlisting}
\par
\medskip

写法拆开看：\texttt{返回类型 fn(\allowbreak{}参数)\allowbreak{} \{ 体 \}}。

\begin{itemize}

\item \texttt{fn} 是关键字；

\item \texttt{(\allowbreak{}int64 x)\allowbreak{}} 跟普通函数的参数一模一样；

\item \texttt{fn} 前面的 \texttt{int64} 是返回类型，可以省略（\texttt{cube} 那份就没写，编译器自己推）；

\item \texttt{\{ … \}} 是函数体。

\end{itemize}

\texttt{let cube =\allowbreak{} …} 让编译器根据右边推出 \texttt{cube} 的类型——省得手写一长串 \texttt{int64 fn(\allowbreak{}int64)\allowbreak{}}。

\section{省略参数类型：靠上下文推}

如果 \texttt{fn} 是要传给某个参数的，而那个参数的类型是已知的，那么连参数类型都能省——编译器按上下文补上：

\par
\medskip
\begin{lstlisting}[style=magel]
int64 apply(int64 fn(int64, int64) f, int64 a, int64 b)
{
    return f(a, b);
}

int64 main()
{
    println("{}", apply(fn(x, y)
    {
        x + y
    }, 3, 4));
    println("{}", apply(fn(x, y)
    {
        x * y
    }, 3, 4));
    return 0;
}
\end{lstlisting}
\par
\medskip

\par
\medskip
\begin{lstlisting}[style=out]
7
12
\end{lstlisting}
\par
\medskip

\texttt{apply} 的第一参数类型是 \texttt{int64 fn(\allowbreak{}int64,\allowbreak{} int64)\allowbreak{}}，所以 \texttt{fn(\allowbreak{}x,\allowbreak{} y)\allowbreak{}} 里的\texttt{x} / \texttt{y} 自动就是 \texttt{int64}。这种「看上下文定类型」的能力，让 \texttt{fn} 写起来非常轻。

\begin{tip}{提示}
但别滥用。参数类型省掉后，读代码的人得自己回头找上下文。类型不是一眼能看出来的地方，就老老实实写上。
\end{tip}

\section{闭包：记住定义处的变量}

\texttt{fn} 最厉害的地方，是它能看见外面的变量：

\par
\medskip
\begin{lstlisting}[style=magel]
int64 main()
{
    int64 base = 10;
    let addBase = fn(int64 x)
    {
        x + base
    };
    base = 100;
    println("{}", addBase(5));
    return 0;
}
\end{lstlisting}
\par
\medskip

\par
\medskip
\begin{lstlisting}[style=out]
105
\end{lstlisting}
\par
\medskip

\texttt{add\allowbreak{}Base(\allowbreak{}5)\allowbreak{}} 读到的是调用那一刻 \texttt{base} 的值 100，所以输出 105。换句话说，闭包里那个 \texttt{base} 不是一份拷贝，而是一条指向外部变量的线。

这条线就叫捕获（capture）：闭包把外部变量"抓"进来，从此你在闭包里读写它，读写的其实是外面那个变量本身。

\begin{tip}{提示}
默认捕获是只读借用。 不写任何修饰时，闭包按 \texttt{const\&} 抓住外部变量：能读、能看到外面的新值，但不能改。想改就得显式声明，见 15.7 节。
\end{tip}

\section{四种捕获方式}

捕获方式写在 \texttt{fn} 关键字后面，一共四种：

{\small\setlength{\tabcolsep}{4pt}
\begin{longtable}{>{\raggedright\arraybackslash}p{\dimexpr 0.294\linewidth-2\tabcolsep\relax}>{\raggedright\arraybackslash}p{\dimexpr 0.185\linewidth-2\tabcolsep\relax}>{\raggedright\arraybackslash}p{\dimexpr 0.152\linewidth-2\tabcolsep\relax}>{\raggedright\arraybackslash}p{\dimexpr 0.196\linewidth-2\tabcolsep\relax}>{\raggedright\arraybackslash}p{\dimexpr 0.152\linewidth-2\tabcolsep\relax}}
\toprule
写法 & 名称 & 闭包里能改？ & 外面那个变量 & 寿命限制 \\
\midrule
\endhead
\texttt{fn const(\allowbreak{}.\allowbreak{}.\allowbreak{}.\allowbreak{})\allowbreak{}}（也是默认） & 只读借用 \texttt{const\&} & 不能改 & 不受影响 & 来源要活过闭包 \\
\texttt{fn mut(\allowbreak{}.\allowbreak{}.\allowbreak{}.\allowbreak{})\allowbreak{}} & 可变借用 \texttt{mut\&} & 能改 & 跟着变 & 来源要活过闭包 \\
\texttt{fn copy(\allowbreak{}.\allowbreak{}.\allowbreak{}.\allowbreak{})\allowbreak{}} & 拥有副本 & 能改（改副本） & 不受影响 & 无，自己有一份 \\
\texttt{fn move(\allowbreak{}.\allowbreak{}.\allowbreak{}.\allowbreak{})\allowbreak{}} & 转移所有权 & 能改（独占） & 被搬走，须重新声明 & 无，自己有一份 \\
\bottomrule
\end{longtable}
}

一句话区分：\texttt{const\&} / \texttt{mut\&} 是"借"，\texttt{copy} / \texttt{move} 是"拿"。 借来的东西要还，所以有寿命限制；拿了就是自己的，想活多久活多久。

\subsection{\texttt{const\&}：默认，只读借用}

不写修饰就是它，写 \texttt{const} 也是它——下面两种写法完全等价：

\par
\medskip
\begin{lstlisting}[style=magel]
int64 main()
{
    int64 n = 1;
    let implicit = int64 fn()
    {
        n
    };
    /* 默认：const& */
    let explicit = int64 fn const()
    {
        n
    };
    /* 写全，同一个意思 */
    n = 2;
    println("{} {}", implicit(), explicit());
    return 0;
}
\end{lstlisting}
\par
\medskip

\par
\medskip
\begin{lstlisting}[style=out]
2 2
\end{lstlisting}
\par
\medskip

两者都跟着外面的新值走（\texttt{n} 改成 2 之后读出来就是 2）。想在闭包里改它，编译器直接拦下：

\par
\medskip
\begin{lstlisting}[style=magel]
int64 main()
{
    int64 base = 10;
    let bump = fn(int64 x)
    {
        base += x;
        return base;
    };
    println("{}", bump(1));
    return 0;
}
\end{lstlisting}
\par
\medskip

\par
\medskip
\begin{lstlisting}[style=out]
magel: 不可变变量（const）不能再赋值：base
\end{lstlisting}
\par
\medskip

\subsection{\texttt{mut\&}：写回来源}

\par
\medskip
\begin{lstlisting}[style=magel]
int64 main()
{
    int64 counter = 0;
    let tick = fn mut()
    {
        counter += 1;
        return counter;
    };
    println("{} {}", tick(), tick());
    println("outside {}", counter);
    return 0;
}
\end{lstlisting}
\par
\medskip

\par
\medskip
\begin{lstlisting}[style=out]
1 2
outside 2
\end{lstlisting}
\par
\medskip

闭包里的 \texttt{counter} 就是外面那个 \texttt{counter}：两次 \texttt{tick(\allowbreak{})\allowbreak{}} 累加到 2，外面读出来也是 2。这是写计数器、累加器最直接的写法。

\subsection{\texttt{copy}：拿一份独立副本}

\par
\medskip
\begin{lstlisting}[style=magel]
int64 main()
{
    int64 n = 5;
    let independent = int64 fn copy(int64 x)
    {
        n += x;
        return n;
    };
    n = 100;
    println("{} {} {}", independent(2), independent(1), n);
    return 0;
}
\end{lstlisting}
\par
\medskip

\par
\medskip
\begin{lstlisting}[style=out]
7 8 100
\end{lstlisting}
\par
\medskip

闭包从 \texttt{n =\allowbreak{} 5} 起家，之后外面把 \texttt{n} 改成 100 跟它没关系：它自己从 5 加到 7、再加到 8，外面始终是 100。

副本对 \texttt{String}、\texttt{List} 这类带资源的值同样成立——闭包内部改的是自己那一份：

\par
\medskip
\begin{lstlisting}[style=magel]
int64 main()
{
    String text = "old";
    let grow = String fn copy()
    {
        text = text + "!";
        return text;
    };
    text = "outside";
    println("{} / {} / {}", grow(), grow(), text);
    return 0;
}
\end{lstlisting}
\par
\medskip

\par
\medskip
\begin{lstlisting}[style=out]
old! / old!! / outside
\end{lstlisting}
\par
\medskip

\par
\medskip
\begin{lstlisting}[style=magel]
int64 main()
{
    List<int64> values;
    values.push(4);
    let bump = int64 fn copy()
    {
        values[0] += 1;
        return values[0];
    };
    println("{} {} {}", bump(), bump(), values[0]);
    return 0;
}
\end{lstlisting}
\par
\medskip

\par
\medskip
\begin{lstlisting}[style=out]
5 6 4
\end{lstlisting}
\par
\medskip

\texttt{copy} 还有个额外好处：即使来源是 \texttt{const} 也能拷——拷出来的是一份新值，动不到原值：

\par
\medskip
\begin{lstlisting}[style=magel]
int64 main()
{
    const int64 fixed = 7;
    let change = int64 fn copy()
    {
        fixed += 1;
        return fixed;
    };
    println("{} / {}", change(), fixed);
    return 0;
}
\end{lstlisting}
\par
\medskip

\par
\medskip
\begin{lstlisting}[style=out]
8 / 7
\end{lstlisting}
\par
\medskip

\subsection{\texttt{move}：把来源整个搬进闭包}

\par
\medskip
\begin{lstlisting}[style=magel]
int64 main()
{
    String moved = "payload";
    let take = String fn move()
    {
        return moved + "-in";
    };
    String moved = "reset";
    println("{} / {}", take(), moved);
    return 0;
}
\end{lstlisting}
\par
\medskip

\par
\medskip
\begin{lstlisting}[style=out]
payload-in / reset
\end{lstlisting}
\par
\medskip

\texttt{move} 之后，原来的 \texttt{moved} 已经被搬空，不能再用——除非重新声明同名变量。上面就是在 \texttt{move} 之后立刻 \texttt{String moved =\allowbreak{} "reset";} 创建一个新的同名绑定。

没有重新声明就接着读写旧绑定，编译器会拦：

\par
\medskip
\begin{lstlisting}[style=magel]
int64 main()
{
    String moved = "payload";
    let take = String fn move()
    {
        return moved;
    };
    println("{}", moved);
    return 0;
}
\end{lstlisting}
\par
\medskip

\par
\medskip
\begin{lstlisting}[style=out]
magel: [MGL1205] 借用来源已移动或销毁，不能继续访问（源对象已经销毁）：moved
\end{lstlisting}
\par
\medskip

\begin{tip}{提示}
\texttt{copy} 还是 \texttt{move}？ 之后还要用原值 \ensuremath{\rightarrow} \texttt{copy}；原值从此归闭包 \ensuremath{\rightarrow} \texttt{move}。对 \texttt{String} / \texttt{List} / 带 \texttt{Drop} 的 struct，\texttt{move} 省一次拷贝，代价是来源作废。
\end{tip}

\subsection{同一个变量既借又 \texttt{move}：看最后一次使用}

创建借用闭包之后，可以把来源移进另一个闭包；关键是旧借用闭包不能在来源作废后继续使用。下面的 \texttt{f(\allowbreak{}1)\allowbreak{}} 会被编译器以 \texttt{MGL1205} 拒绝，不会进入 VM 执行：

\par
\medskip
\begin{lstlisting}[style=magel]
void main()
{
    int64 n = 1;
    let f = int64 fn mut(int64 x)
    {
        n += x;
        return n;
    };
    let g = int64 fn move(int64 x)
    {
        return x * n;
    };
    println("f={}", f(1));
    println("g={}", g(1));
}
\end{lstlisting}
\par
\medskip

问题出在第三步：\texttt{f} 借用 \texttt{n}，\texttt{g} 随后把 \texttt{n} 搬走，最后却又调用 \texttt{f}。闭包的副本仍然借用原来的来源，复制 \texttt{f} 也不能绕开这个检查。显式销毁或释放来源之后再使用借用闭包，同样会被拒绝。

如果先调用完 \texttt{f}，再创建 \texttt{g}，并且此后不再使用 \texttt{f}，这个顺序可以通过：

\par
\medskip
\begin{lstlisting}[style=magel]
int64 main()
{
    int64 n = 1;
    let f = int64 fn mut(int64 x)
    {
        n += x;
        return n;
    };
    println("{}", f(1));
    let g = int64 fn move(int64 x)
    {
        return x * n;
    };
    println("{}", g(3));
    return 0;
}
\end{lstlisting}
\par
\medskip

\par
\medskip
\begin{lstlisting}[style=out]
2
6
\end{lstlisting}
\par
\medskip

两个闭包需要交替使用时，可以让其中一个拥有快照：

\par
\medskip
\begin{lstlisting}[style=magel]
int64 main()
{
    int64 n = 1;
    let f = fn mut(x) { n = x + n };
    let g = fn copy(x) { x * n };
    println("{} {} {}", f(1), g(1), n);
    return 0;
}
\end{lstlisting}
\par
\medskip

\par
\medskip
\begin{lstlisting}[style=out]
2 1 2
\end{lstlisting}
\par
\medskip

\texttt{g} 用 \texttt{fn copy} 拿走的是创建闭包那一刻的快照（\texttt{n =\allowbreak{} 1}），\texttt{n} 不会再被搬空，\texttt{f} 那条线始终有效。

\section{捕获列表：逐个变量指定方式}

\texttt{fn mut(\allowbreak{}.\allowbreak{}.\allowbreak{}.\allowbreak{})\allowbreak{}} 是"所有外部变量都按 \texttt{mut\&} 抓"。想给不同变量不同待遇，就在 \texttt{fn} 后面加一对方括号：

\par
\medskip
\begin{lstlisting}[style=magel]
int64 main()
{
    int64 a = 2;
    int64 b = 3;
    int64 c = 4;
    int64 d = 5;
    let mixed = int64 fn copy[a] mut[b] const[c] move[d](int64 x)
    {
        a += x;
        b += x;
        d += x;
        return a + b + c + d;
    };
    a = 100;
    c = 10;
    d = 200;
    println("{} | a={} b={} c={} d={}", mixed(1), a, b, c, d);
    println("{} | b={}", mixed(1), b);
    return 0;
}
\end{lstlisting}
\par
\medskip

\par
\medskip
\begin{lstlisting}[style=out]
23 | a=100 b=4 c=10 d=200
26 | b=5
\end{lstlisting}
\par
\medskip

逐个看：\texttt{a} 拷了一份（外面改 100 不影响它），\texttt{b} 借用（闭包加 1 后外面就是 4），\texttt{c} 只读借用（不能改，但能看到新值 10），\texttt{d} 被搬进闭包（外面重新赋成 200 是另一回事）。

其余变量的默认捕获方式。 列表后面还能指定未单独列出的变量采用什么方式捕获。

\par
\medskip
\begin{lstlisting}[style=magel]
fn copy[a, b] mut(int64 x)
{
    ...
}
/*     ^^^^^^^^^^ ^^^                         */
/*     a、b 按 copy 抓，其余变量按 mut& 抓     */
\end{lstlisting}
\par
\medskip

\par
\medskip
\begin{lstlisting}[style=magel]
int64 main()
{
    int64 left = 1;
    int64 right = 2;
    int64 shared = 3;
    let f = int64 fn copy[left, right] mut(int64 x)
    {
        left += x;
        right += x;
        shared += x;
        return left + right + shared;
    };
    println("{} | left={} right={} shared={}", f(1), left, right, shared);
    return 0;
}
\end{lstlisting}
\par
\medskip

\par
\medskip
\begin{lstlisting}[style=out]
9 | left=1 right=2 shared=4
\end{lstlisting}
\par
\medskip

\texttt{left} / \texttt{right} 是副本（外面不动），\texttt{shared} 没列出、按默认 \texttt{mut\&} 借用，所以闭包加完之后外面读到的 \texttt{shared} 变成了 4。

不写默认方式时，没列出的变量一律按 \texttt{const\&} 抓——和全局默认一致。默认方式只能写一次，写两个会报错：

\par
\medskip
\begin{lstlisting}[style=magel]
int64 main()
{
    int64 a = 1;
    let f = int64 fn copy[a] mut const()
    {
        a
    };
    println("{}", f());
    return 0;
}
\end{lstlisting}
\par
\medskip

\par
\medskip
\begin{lstlisting}[style=out]
magel: lambda：lambda 默认捕获方式只能声明一次：mut 与 const
\end{lstlisting}
\par
\medskip

方括号也可以是空的——\texttt{fn copy[\allowbreak{}]\allowbreak{}(\allowbreak{}.\allowbreak{}.\allowbreak{}.\allowbreak{})\allowbreak{}} 表示"一个都不特殊对待"，其余变量仍按默认 \texttt{const\&}：

\par
\medskip
\begin{lstlisting}[style=magel]
int64 main()
{
    int64 bias = 9;
    let empty = int64 fn copy[]()
    {
        bias
    };
    bias = 10;
    println("{}", empty());
    return 0;
}
\end{lstlisting}
\par
\medskip

\par
\medskip
\begin{lstlisting}[style=out]
10
\end{lstlisting}
\par
\medskip

传给函数参数的 \texttt{fn} 同样能带捕获列表：

\par
\medskip
\begin{lstlisting}[style=magel]
int64 apply(int64 fn(int64) f, int64 x)
{
    return f(x);
}
int64 main()
{
    int64 bias = 9;
    println("{} / {}", apply(fn copy[bias](x)
    {
        bias += x;
        return bias;
    }, 2), bias);
    return 0;
}
\end{lstlisting}
\par
\medskip

\par
\medskip
\begin{lstlisting}[style=out]
11 / 9
\end{lstlisting}
\par
\medskip

\section{复制闭包：拥有状态会复制，借用仍然共享}

闭包是值，能赋给另一个变量。那复制过去的到底是什么？

答案是：拥有的东西各复制一份，借的东西还是同一条线。

\texttt{copy} 捕获的闭包复制后，两份副本互不相干：

\par
\medskip
\begin{lstlisting}[style=magel]
int64 main()
{
    int64 n = 5;
    let a = int64 fn copy(int64 x)
    {
        n += x;
        return n;
    };
    let b = a;
    println("b={} a={} b={}", b(3), a(1), b(1));
    return 0;
}
\end{lstlisting}
\par
\medskip

\par
\medskip
\begin{lstlisting}[style=out]
b=8 a=6 b=9
\end{lstlisting}
\par
\medskip

\texttt{a}、\texttt{b} 各有一份从 5 开始的副本：\texttt{b} 先加到 8，\texttt{a} 加到 6，\texttt{b} 再加到 9——三笔独立的账。

而 \texttt{mut\&} 捕获的闭包复制后，两个闭包仍指向同一个来源：

\par
\medskip
\begin{lstlisting}[style=magel]
int64 main()
{
    int64 n = 0;
    let a = int64 fn mut(int64 x)
    {
        n += x;
        return n;
    };
    let b = a;
    println("a={} b={} outside={}", a(1), b(1), n);
    return 0;
}
\end{lstlisting}
\par
\medskip

\par
\medskip
\begin{lstlisting}[style=out]
a=1 b=2 outside=2
\end{lstlisting}
\par
\medskip

\texttt{a} 加完 \texttt{n} 是 1，\texttt{b} 接着加变成 2，外面读也是 2——同一个 \texttt{n}，同一条线。

\section{生命周期：借来的，不能活过来源}

\texttt{const\&} / \texttt{mut\&} 是借用，就有一条铁律：

\begin{tip}{提示}
闭包（以及它的任何副本）不能比它借的那个变量活得更久。
\end{tip}

最容易踩的地方是把闭包返回出去：

\par
\medskip
\begin{lstlisting}[style=magel]
int64 fn() counter()
{
    int64 n = 4;
    return int64 fn mut()
    {
        n += 1;
        return n;
    };
}
int64 main()
{
    let c = counter();
    println("{}", c());
    return 0;
}
\end{lstlisting}
\par
\medskip

\par
\medskip
\begin{lstlisting}[style=out]
magel: [MGL1204] 返回值中的引用不能超过局部来源的生命周期
\end{lstlisting}
\par
\medskip

\texttt{n} 是 \texttt{counter} 的局部变量，函数一返回它就没了；闭包要是带走了它的引用，调用时访问的就是一块已经死掉的内存。编译器不放行（报错里还会附上那个来源的名字）。

换成 \texttt{copy} 或 \texttt{move}，闭包自己拥有一份，就能安全返回：

\par
\medskip
\begin{lstlisting}[style=magel]
int64 fn() counter()
{
    int64 n = 4;
    return int64 fn copy()
    {
        n += 1;
        return n;
    };
}
int64 main()
{
    let c = counter();
    println("{} {}", c(), c());
    return 0;
}
\end{lstlisting}
\par
\medskip

\par
\medskip
\begin{lstlisting}[style=out]
5 6
\end{lstlisting}
\par
\medskip

\texttt{move} 同理——把 \texttt{n} 搬进闭包，闭包从此自带状态：

\par
\medskip
\begin{lstlisting}[style=magel]
int64 fn() counter()
{
    int64 n = 8;
    return int64 fn move()
    {
        n += 1;
        return n;
    };
}
int64 main()
{
    let c = counter();
    println("{} {}", c(), c());
    return 0;
}
\end{lstlisting}
\par
\medskip

\par
\medskip
\begin{lstlisting}[style=out]
9 10
\end{lstlisting}
\par
\medskip

判断口诀：闭包要返回 \ensuremath{\rightarrow} 捕获只能是 \texttt{copy} / \texttt{move}，或者是借"调用方"的东西。

最后半句正是 15.14 节"闭包当工厂"的做法：借用调用方传进来的指针，那个变量活在外面，闭包带回去才安全。

\section{带 \texttt{Drop} 的类型：谁负责销毁}

如果捕获的是带 \texttt{Drop} 的类型（有 \texttt{void drop(\allowbreak{})\allowbreak{}} 的 struct），所有权规则就直接体现为什么时候调用 \texttt{drop(\allowbreak{})\allowbreak{}}。

\texttt{copy} 捕获 \ensuremath{\rightarrow} 外面那个值和闭包里的副本是两个对象，各销毁一次：

\par
\medskip
\begin{lstlisting}[style=magel]
int64 drops = 0;
struct Record: Drop
{
    String text;
    List<int64> values;
    void drop()
    {
        drops += 1;
    }
};
int64 main()
{
    Record record;
    record.text = "owned";
    record.values.push(1);
    {
        let own = int64 fn copy()
        {
            record.values[0] += 1;
            return record.values[0];
        };
        println("own={}", own());
        {
            let second = own;
            println("second={} drops={}", second(), drops);
        }
        println("drops={}", drops);
    }
    println("drops={} outside={}", drops, record.values[0]);
    return 0;
}
\end{lstlisting}
\par
\medskip

\par
\medskip
\begin{lstlisting}[style=out]
own=2
second=3 drops=0
drops=1
drops=2 outside=1
\end{lstlisting}
\par
\medskip

\texttt{second} 出作用域销毁一份副本（\texttt{drops=\allowbreak{}1}），\texttt{own} 出作用域再销毁一份（\texttt{drops=\allowbreak{}2}）。全程外面 \texttt{record.\allowbreak{}values[\allowbreak{}0]\allowbreak{}} 一直是 1——副本里加多少都影响不到它。

\texttt{move} 捕获 \ensuremath{\rightarrow} 对象被搬进闭包，只有闭包销毁时才 \texttt{drop(\allowbreak{})\allowbreak{}}：

\par
\medskip
\begin{lstlisting}[style=magel]
int64 drops = 0;
struct Record: Drop
{
    String text;
    List<int64> values;
    void drop()
    {
        drops += 1;
    }
};
int64 main()
{
    Record record;
    record.text = "moved";
    record.values.push(10);
    {
        let own = int64 fn move()
        {
            record.values[0] += 1;
            return record.values[0];
        };
        println("own={} drops={}", own(), drops);
    }
    println("drops={}", drops);
    return 0;
}
\end{lstlisting}
\par
\medskip

\par
\medskip
\begin{lstlisting}[style=out]
own=11 drops=0
drops=1
\end{lstlisting}
\par
\medskip

闭包出作用域 \ensuremath{\rightarrow} 它拥有的那个 \texttt{Record} 被销毁 \ensuremath{\rightarrow} \texttt{drops} 从 0 变 1，全程只有一次。

\section{捕获方式速查}

{\small\setlength{\tabcolsep}{4pt}
\begin{longtable}{>{\raggedright\arraybackslash}p{\dimexpr 0.394\linewidth-2\tabcolsep\relax}>{\raggedright\arraybackslash}p{\dimexpr 0.125\linewidth-2\tabcolsep\relax}>{\raggedright\arraybackslash}p{\dimexpr 0.115\linewidth-2\tabcolsep\relax}>{\raggedright\arraybackslash}p{\dimexpr 0.115\linewidth-2\tabcolsep\relax}>{\raggedright\arraybackslash}p{\dimexpr 0.231\linewidth-2\tabcolsep\relax}}
\toprule
写法 & 闭包内能改？ & 外面跟着变？ & 来源还能用？ & 能返回出去？ \\
\midrule
\endhead
\texttt{fn(\allowbreak{}.\allowbreak{}.\allowbreak{}.\allowbreak{})\allowbreak{}} / \texttt{fn const(\allowbreak{}.\allowbreak{}.\allowbreak{}.\allowbreak{})\allowbreak{}} & 否 & — & 能 & 借用来源寿命足够时能 \\
\texttt{fn mut(\allowbreak{}.\allowbreak{}.\allowbreak{}.\allowbreak{})\allowbreak{}} & 是 & 是 & 能 & 借用来源寿命足够时能 \\
\texttt{fn copy(\allowbreak{}.\allowbreak{}.\allowbreak{}.\allowbreak{})\allowbreak{}} & 是 & 否 & 能 & 不含短寿命借用时能 \\
\texttt{fn move(\allowbreak{}.\allowbreak{}.\allowbreak{}.\allowbreak{})\allowbreak{}} & 是 & 否 & 需重新声明 & 不含短寿命借用时能 \\
\texttt{fn copy[\allowbreak{}a]\allowbreak{} mut[\allowbreak{}b]\allowbreak{} const[\allowbreak{}c]\allowbreak{} move[\allowbreak{}d]\allowbreak{}(\allowbreak{}.\allowbreak{}.\allowbreak{}.\allowbreak{})\allowbreak{}} & 各按各的 & 各按各的 & 各按各的 & 所有借用来源寿命足够时能 \\
\texttt{fn copy[\allowbreak{}a,\allowbreak{} b]\allowbreak{} mut(\allowbreak{}.\allowbreak{}.\allowbreak{}.\allowbreak{})\allowbreak{}} & 其余按 \texttt{mut\&} & — & — & 取决于借用来源寿命 \\
\bottomrule
\end{longtable}
}

容易忘的几条：

\begin{itemize}

\item 默认就是 \texttt{const\&}——不写 \texttt{mut} 就改不了外面的变量。

\item \texttt{mut} 改的是外面那个，不是副本；想各算各的用 \texttt{copy}。

\item \texttt{move} 之后来源作废，须重新声明新绑定，普通赋值不能恢复。

\item 借用本次调用局部变量的闭包不能返回；需要自有状态时用 \texttt{copy} / \texttt{move}。借用调用方提供的对象仍须满足来源寿命。

\item 借用闭包的最后一次使用要在来源作废之前：先用完再 \texttt{move} 可以通过；\texttt{move} 后再使用旧借用闭包会报 \texttt{MGL1205}（见 15.7.5）。

\item 复制闭包：拥有的状态各复制一份，借用的线仍指向同一来源。

\end{itemize}

\section{一个小坑：表达式体 vs 语句体}

\texttt{fn} 的体有两种写法，差别就一个分号：

\par
\medskip
\begin{lstlisting}[style=magel]
let f = fn(int64 x)
{
    x * 2
};
// 表达式体：没有分号
let g = fn(int64 x)
{
    return x * 2;
};
// 语句体：有分号 + return
\end{lstlisting}
\par
\medskip

表达式体是「一个表达式」，直接就是返回值——末尾不能带分号。一旦你手滑写成 \texttt{\{ x * 3; \}}（带分号），它就成了一「条语句」，编译器会要求你补 \texttt{return}：

\par
\medskip
\begin{lstlisting}[style=magel]
int64 main()
{
    let g = int64 fn(int64 x)
    {
        x * 3;
    };
    println("{}", g(5));
    return 0;
}
\end{lstlisting}
\par
\medskip

\par
\medskip
\begin{lstlisting}[style=out]
magel: 返回类型为 int64 的函数 __magel_lambda_… 不是所有路径都有 return
\end{lstlisting}
\par
\medskip

记住一句话：体里想要「最后一行的值」当返回值 \ensuremath{\rightarrow} 别加分号；想用 \texttt{return} \ensuremath{\rightarrow} 加分号。

\section{闭包当工厂：返回一个函数}

15.10 节说过：借用的闭包不能返回。想让工厂函数产出一个闭包，捕获来源就得由调用方提供——让调用方把变量的地址传进来，那个变量活在外面，闭包把它带回去才安全：

\par
\medskip
\begin{lstlisting}[style=magel]
int64 fn(int64) makeAdder(int64 mut& amount)
{
    return fn(int64 value)
    {
        value + amount
    };
}
int64 main()
{
    int64 one = 1;
    int64 ten = 10;
    let plus1 = makeAdder(&one);
    let plus10 = makeAdder(&ten);
    println("{} {}", plus1(5), plus10(5));
    return 0;
}
\end{lstlisting}
\par
\medskip

\par
\medskip
\begin{lstlisting}[style=out]
6 15
\end{lstlisting}
\par
\medskip

\texttt{make\allowbreak{}Adder(\allowbreak{}\&one)\allowbreak{}} 和 \texttt{make\allowbreak{}Adder(\allowbreak{}\&ten)\allowbreak{}} 返回的 \texttt{plus1} / \texttt{plus10}，各自借着\texttt{one} / \texttt{ten}：只要这两个变量还在，两个闭包就一直能用；它们要是先没了，闭包也就跟着失效——这正是 15.10 那条铁律的另一面。

函数值本身可以复制；复制闭包时到底复制了什么，见 15.9 节。

\subsection*{移动后访问诊断}

下面的闭包将 a 的所有权搬进自身，创建之后读取源变量 a 会报 MGL1205：

\par
\medskip
\begin{lstlisting}[style=magel]
int64 main()
{
    String a = "1";
    String fn() b = fn move()
    {
        a
    };
    println("{}", a);
    return 0;
}
\end{lstlisting}
\par
\medskip

如果还要保留 a，就按值复制捕获；如果使用 move，之后通过 b() 读取闭包拥有的值。

\par
\medskip
\begin{lstlisting}[style=magel]
int64 main()
{
    String a = "1";
    String fn() b = fn copy()
    {
        a
    };
    println("{} {}", a, b());
    return 0;
}
\end{lstlisting}
\par
\medskip

\par
\medskip
\begin{lstlisting}[style=out]
1 1
\end{lstlisting}
\par
\medskip

编辑器对捕获位置的 move/copy/const/mut 显示对应说明，不套用普通 \texttt{move<\allowbreak{}T>\allowbreak{}} 函数签名；诊断红线标到后续读取的变量。

\section{语句级具名闭包}

如果你希望给局部闭包直接起名字，可以写成 \texttt{返回类型 名字(\allowbreak{}参数)\allowbreak{} \{ .\allowbreak{}.\allowbreak{}.\allowbreak{} \}}。它仍是函数值，只在声明语句位置识别，不能塞进另一个表达式中间。返回类型是检查函数体的类型锚点。

\par
\medskip
\begin{lstlisting}[style=magel]
int64 fn() makeNamedCounter()
{
    int64 current = 0;
    int64 next() move
    {
        current++;
        return current;
    }
    return next;
}

int64 main()
{
    int64 base = 10;
    int64 read()
    {
        return base;
    }
    int64 increment() mut
    {
        base++;
        return base;
    }
    int64 isolated() copy
    {
        base += 100;
        return base;
    }
    int64 initial = read();
    int64 changed = increment();
    int64 snapshot = isolated();
    println("{} {} {} {}", initial, changed, snapshot, base);

    let counter = makeNamedCounter();
    int64 first = counter();
    int64 second = counter();
    println("{} {}", first, second);

    int64 callback(int64 value);
    callback = int64 fn(int64 value)
    {
        return value + 1;
    };
    println("{}", callback(5));
    return 0;
}
\end{lstlisting}
\par
\medskip

\par
\medskip
\begin{lstlisting}[style=out]
10 11 110 11
1 2
6
\end{lstlisting}
\par
\medskip

默认捕获方式是 \texttt{const\&}，因此 \texttt{read} 能看到来源后续的变化，却不能修改来源。函数体前显式写 \texttt{mut} 是可写借用，\texttt{copy} 拥有独立状态，\texttt{move} 则接管捕获值的所有权。捕获状态在声明初始化时交给函数名绑定；普通后续赋值仍遵守值语义。

局部 \texttt{int64 callback(\allowbreak{}int64 value)\allowbreak{};} 只声明一个有确定签名的函数变量，后面再赋函数值；它与模块层的宿主函数声明不是同一个语义入口。借用捕获的闭包不能逃出来源的生命周期。工厂函数中的 \texttt{next} 使用 move 拥有状态，所以可安全返回；把它改成默认借用并返回会报错。

\section{多层高阶闭包：函数也可以接收函数}

\texttt{R fn(\allowbreak{}A)\allowbreak{}} 表示“接收 \texttt{A}、返回 \texttt{R} 的函数值”。这里的 \texttt{A} 也可以是函数类型，因此 \texttt{fn} 可以继续嵌套。下面保留五层嵌套的完整写法，并通过参数名和换行区分两个实参：

\par
\medskip
\begin{lstlisting}[style=magel]
int64 hello(
    int64 fn(int64 fn(int64 fn(int64 fn(int64 fn(int64))))) pipeline,
    int64 fn(int64) finish)
{
    return finish(pipeline(level4));
}
\end{lstlisting}
\par
\medskip

这里是具名函数声明；有函数体时，模块层声明在右大括号结束，后面不加分号。形参要有名字，函数体才能调用它们。\texttt{fn} 后的小括号属于函数类型，不是一次函数调用。

由内向外阅读第一项形参的类型：

\begin{center}
\begin{tabular}{llll}
\toprule
层次 & 接收什么 & 返回什么 & 下方程序中的值 \\
\midrule
第一层 & \texttt{int64} & \texttt{int64} & \texttt{level1} \\
第二层 & 第一层函数 & \texttt{int64} & \texttt{level2} \\
第三层 & 第二层函数 & \texttt{int64} & \texttt{level3} \\
第四层 & 第三层函数 & \texttt{int64} & \texttt{level4} \\
第五层 & 第四层函数 & \texttt{int64} & \texttt{level5} \\
\bottomrule
\end{tabular}
\end{center}

这些层次表示参数类型的嵌套深度，不表示连续调用五次后才得到一个返回函数；表中每一层的返回类型都是 \texttt{int64}。\texttt{hello} 的第一项实参是第五层函数，第二项实参是普通的一元闭包。

完整程序如下。最后的 \texttt{finish} 捕获了 \texttt{main} 中的 \texttt{factor}，因此这同时展示了高阶函数参数与带状态的闭包：

\par
\medskip
\begin{lstlisting}[style=magel]
int64 level1(int64 value)
{
    return value + 1;
}

int64 level2(int64 fn(int64) callback)
{
    return callback(10);
}

int64 level3(int64 fn(int64 fn(int64)) callback)
{
    return callback(level1);
}

int64 level4(int64 fn(int64 fn(int64 fn(int64))) callback)
{
    return callback(level2);
}

int64 level5(int64 fn(int64 fn(int64 fn(int64 fn(int64)))) callback)
{
    return callback(level3);
}

int64 hello(
    int64 fn(int64 fn(int64 fn(int64 fn(int64 fn(int64))))) pipeline,
    int64 fn(int64) finish)
{
    return finish(pipeline(level4));
}

int64 main()
{
    int64 factor = 2;
    int64 finish(int64 value)
    {
        return value * factor;
    }
    println("{}", hello(level5, finish));
    return 0;
}
\end{lstlisting}
\par
\medskip

\par
\medskip
\begin{lstlisting}[style=out]
22
\end{lstlisting}
\par
\medskip

调用顺序是 \texttt{pipeline(\allowbreak{}level4)\allowbreak{}} \ensuremath{\rightarrow} \texttt{level4(\allowbreak{}level3)\allowbreak{}} \ensuremath{\rightarrow} \texttt{level3(\allowbreak{}level2)\allowbreak{}} \ensuremath{\rightarrow} \texttt{level2(\allowbreak{}level1)\allowbreak{}} \ensuremath{\rightarrow} \texttt{level1(\allowbreak{}10)\allowbreak{}}，先得到 \texttt{11}；随后 \texttt{finish(\allowbreak{}11)\allowbreak{}} 读取捕获的 \texttt{factor}，得到 \texttt{22}。具名函数和闭包都按各自的完整函数类型检查，某一层参数类型不匹配就不能传入。

\texttt{finish} 默认以 \texttt{const\&} 借用 \texttt{factor}。这次调用在 \texttt{main} 的作用域内同步完成，借用始终有效。若要把这样的闭包返回或存到更长寿命的对象中，仍须满足前面介绍的捕获与生命周期规则；嵌套 \texttt{fn} 类型不会延长被借用变量的寿命。

\section{组合闭包：接收函数，再返回函数}

高阶函数还可以把几个操作组合成新的函数值。下面的 \texttt{combine} 接收两个 \texttt{int64 fn(\allowbreak{}int64)\allowbreak{}}，返回的类型也是 \texttt{int64 fn(\allowbreak{}int64)\allowbreak{}}：

\par
\medskip
\begin{lstlisting}[style=magel]
let combine = (int64 fn(int64)) fn(int64 fn(int64) a, int64 fn(int64) b)
{
    int64 fn copy(int64 x)
    {
        a(b(x))
    }
};
\end{lstlisting}
\par
\medskip

从左到右拆开这个声明：括号里的 \texttt{int64 fn(\allowbreak{}int64)\allowbreak{}} 是返回类型，接下来的 \texttt{fn(\allowbreak{}.\allowbreak{}.\allowbreak{}.\allowbreak{})\allowbreak{}} 才是创建组合器的匿名函数。它的函数体每次创建一个新闭包，先调用 \texttt{b}，再把结果交给 \texttt{a}。内层的 \texttt{a(\allowbreak{}b(\allowbreak{}x)\allowbreak{})\allowbreak{}} 与外层的闭包表达式都处在块末尾且不加分号，分别成为两层函数的返回值。

返回函数类型外面的括号用于分组，不会增加一层函数类型；也可以直接写 \texttt{int64 fn(\allowbreak{}int64)\allowbreak{} fn(\allowbreak{}.\allowbreak{}.\allowbreak{}.\allowbreak{})\allowbreak{}}。参数列表最后一个参数后面不加逗号。与上一节的具名函数定义不同，这里是变量初始化语句，最外层右大括号后面需要分号。

完整程序如下：

\par
\medskip
\begin{lstlisting}[style=magel]
int64 main()
{
    let combine = (int64 fn(int64)) fn(int64 fn(int64) a, int64 fn(int64) b)
    {
        int64 fn copy(int64 x)
        {
            a(b(x))
        }
    };

    let f1 = int64 fn(int64 x) { x + 1 };
    let f2 = int64 fn(int64 x) { x * 2 };
    let c1 = combine(f1, f2);
    let c2 = combine(f2, f1);
    println("{} {}", c1(5), c2(5));

    let c3 = combine(fn(x) { x * 5 }, fn(x) { x - 1 });
    println("{}", c3(5));
    return 0;
}
\end{lstlisting}
\par
\medskip

\par
\medskip
\begin{lstlisting}[style=out]
11 12
20
\end{lstlisting}
\par
\medskip

\texttt{c1(\allowbreak{}5)\allowbreak{}} 先执行 \texttt{f2(\allowbreak{}5)\allowbreak{}} 得到 \texttt{10}，再执行 \texttt{f1(\allowbreak{}10)\allowbreak{}} 得到 \texttt{11}；\texttt{c2(\allowbreak{}5)\allowbreak{}} 先执行 \texttt{f1(\allowbreak{}5)\allowbreak{}} 得到 \texttt{6}，再执行 \texttt{f2(\allowbreak{}6)\allowbreak{}} 得到 \texttt{12}。两次调用各自返回一个拥有 \texttt{a}、\texttt{b} 副本的组合闭包，交换实参会改变组合顺序。

内层函数不能默认借用本次调用的局部形参 \texttt{a}、\texttt{b} 后直接返回；这些形参会随 \texttt{combine} 返回而离开作用域。\texttt{fn copy} 使返回闭包拥有这两个函数值的副本。本例的 \texttt{f1}、\texttt{f2} 都没有捕获外部变量；如果它们含有借用，复制函数值不会把被借用对象变成自有状态，其来源仍必须活得足够久。

\texttt{c3} 直接把两个匿名函数作为实参传入。目标形参 \texttt{int64 fn(\allowbreak{}int64)\allowbreak{}} 为 \texttt{fn(\allowbreak{}x)\allowbreak{}} 提供参数类型与返回类型的上下文，因此两个 \texttt{x} 都推断为 \texttt{int64}，不必重复写类型。组合顺序仍然是先执行第二项实参，再执行第一项实参：\texttt{(\allowbreak{}5 - 1)\allowbreak{} * 5 =\allowbreak{} 20}。这些匿名函数没有捕获状态，传入后同样由组合闭包持有其副本。

\section{泛型闭包：同一份状态，多种参数类型}

闭包也能声明类型参数。\texttt{fn<\allowbreak{}T>\allowbreak{}(\allowbreak{}T value)\allowbreak{}} 中的 \texttt{T} 在调用时确定，通常可以从实参推导；需要指定时写 \texttt{f<\allowbreak{}int64>\allowbreak{}(\allowbreak{}3)\allowbreak{}}。具体实现体决定返回类型。

特化沿用主定义的名称：\texttt{f<\allowbreak{}int64>\allowbreak{} =\allowbreak{} fn(\allowbreak{}int64 value)\allowbreak{} \{ .\allowbreak{}.\allowbreak{}.\allowbreak{} \};} 是完整特化，\texttt{f<\allowbreak{}T,\allowbreak{} int64>\allowbreak{} =\allowbreak{} fn(\allowbreak{}T first,\allowbreak{} int64 second)\allowbreak{} \{ .\allowbreak{}.\allowbreak{}.\allowbreak{} \};} 是偏特化。编译器选择唯一最具体的匹配；重复定义和无法区分的重叠模式会报错。它们是编译期声明，不是在运行时给闭包重新赋值。

泛型闭包可以直接成为结构体成员。普通闭包和泛型闭包也都可以传给泛型类型参数 \texttt{F}；\texttt{apply} 里的 \texttt{F} 由实参推导，不必写出一个固定的函数指针签名。

\par
\medskip
\begin{lstlisting}[style=magel]
struct Mapper
{
    let transform = fn<T>(T value)
    {
        value
    };
    transform<int64> = fn(int64 value)
    {
        value + 1
    };
};

int64 apply<F>(F operation, int64 value)
{
    return operation(value);
}

int64 main()
{
    let identity = fn<T>(T value)
    {
        value
    };
    Mapper mapper;
    println("{} {} {}", mapper.transform(2), mapper.transform<String>("member"), apply(identity, 4));
    return 0;
}
\end{lstlisting}
\par
\medskip

\par
\medskip
\begin{lstlisting}[style=out]
3 member 4
\end{lstlisting}
\par
\medskip

闭包作为 \texttt{F} 传递后，仍保留捕获环境，也能放进 \texttt{Box<\allowbreak{}F>\allowbreak{}} 这样的泛型聚合。借用来源的寿命要求不会因为经过泛型而消失。

不同类型的调用共享同一个闭包的状态。下面的 \texttt{count} 在创建 \texttt{next} 时复制一次；\texttt{bool}、\texttt{int64} 和 \texttt{String} 调用都修改这份副本，完整特化也不另建环境。

\par
\medskip
\begin{lstlisting}[style=magel]
int64 main()
{
    int64 count = 0;
    let next = fn<T> copy[count](T value)
    {
        count += 1;
        return count;
    };
    next<int64> = fn(int64 value)
    {
        count += 10;
        return count;
    };
    println("{} {} {} {}", next(true), next(2), next("text"), count);
    return 0;
}
\end{lstlisting}
\par
\medskip

\par
\medskip
\begin{lstlisting}[style=out]
1 11 12 0
\end{lstlisting}
\par
\medskip

特化沿用主定义的捕获方式；要获得另一份独立状态，复制整个闭包。默认借用、\texttt{mut}、\texttt{copy}、\texttt{move} 仍按本章前面的规则处理，来源作废后使用借用闭包仍会被拒绝。

\section{小结}

\begin{center}
\begin{tabular}{ll}
\toprule
你会了 & 关键点 \\
\midrule
函数类型 & \texttt{R fn(\allowbreak{}A,\allowbreak{} B)\allowbreak{}}，和别的类型一样能当变量/字段/参数 \\
函数当值 & \texttt{int64 fn(\allowbreak{}int64,\allowbreak{} int64)\allowbreak{} op =\allowbreak{} add;}，可重新指向 \\
传函数 & 把函数名或 \texttt{fn} 当实参传，实现回调 \\
\texttt{fn} & \texttt{返回类型 fn(\allowbreak{}参数)\allowbreak{} \{ 体 \}}，返回类型可省 \\
上下文推断 & 目标类型已知时，参数类型也能省 \\
闭包 & 记住定义处的变量；默认 \texttt{const\&} 只读借用 \\
四种捕获 & \texttt{const\&} 只读、\texttt{mut\&} 写回、\texttt{copy} 独立副本、\texttt{move} 转移所有权 \\
捕获列表 & \texttt{fn copy[\allowbreak{}a]\allowbreak{} mut[\allowbreak{}b]\allowbreak{} const[\allowbreak{}c]\allowbreak{} move[\allowbreak{}d]\allowbreak{}(\allowbreak{}.\allowbreak{}.\allowbreak{}.\allowbreak{})\allowbreak{}}；列表后可加默认方式 \\
复制闭包 & 拥有状态各复制一份，借用仍指向同一来源 \\
生命周期 & 借用的闭包（及其副本）不能比来源活得更久 \\
返回闭包 & 只有 \texttt{copy} / \texttt{move}、或借调用方的指针，才能安全带出函数 \\
\texttt{Drop} & \texttt{copy} 两份各销毁一次，\texttt{move} 只在闭包销毁时销毁一次 \\
体例坑 & 表达式体别加分号；用 \texttt{return} 就要加分号 \\
工厂 & 返回闭包，配不同的「料」造出不同的函数 \\
\bottomrule
\end{tabular}
\end{center}

\section{下一步}

\texttt{filter}、\texttt{map} 这类「把函数当参数」的操作，写起来越来越顺手了。但你自己写 \texttt{filter} 还得手搓 \texttt{List}、手搓 \texttt{for-in}——一不小心就写出一堆模板代码。

第 16 章讲错误处理，先把「出错了怎么办」这件事说清楚；之后第六部分的迭代器，会把 \texttt{map} / \texttt{filter} 这些套路做成一条能直接用的链。

翻到\hyperref[ch:16]{第 16 章 错误处理}。

\section{边界与验证}

拥有捕获与借用捕获影响返回和保存闭包的合法性。可变闭包、闭包副本和多次调用的状态变化分别验证；返回闭包不等于允许返回对局部对象的借用。

\section{习题}

前两题用于理解和基本实现，第三题检查边界，最后一题为综合挑战。先独立完成，再核对参考解答。

\begin{enumerate}

\item 默认捕获、mut、copy 和 move 分别表达什么关系？

\item 用 copy 捕获构造可返回的加法闭包。

\item 复制一个有状态闭包，再分别调用并比较状态。

\item 挑战：将一个捕获局部变量而逃逸的反例修复为合法版本。

\end{enumerate}

\subsection*{参考解答与验收要点}

\begin{enumerate}

\item 默认读取来源，mut 允许借用写回，copy 拥有独立副本，move 转移拥有关系；仍需遵守各自类型和生命周期要求。

\item 让返回结果拥有所需状态，离开构造函数后调用仍有效；不要以延长局部变量存活作为解释。

\item 记录两份闭包各自的返回序列；与显式共享对象的捕获情形分开。

\item 先证明反例因借用来源失效被拒绝，再根据需要选择拥有副本或转移，验证调用结果。

\end{enumerate}

\section{番外：孔乙己与闭包}

改写自鲁迅《孔乙己》，正文采用本书读者提供的改编稿。人物对白中的判断不代替本章的语言规则。

鲁镇的码栈的格局，是和别处不同的：都是当街一个曲尺形的大柜台，柜里面预备着服务器，可以随时跑代码。做工的人，傍午傍晚散了工，每每花四文铜钱，开一个沙箱——这是二十多年前的事，现在每个沙箱要涨到十文——靠柜外站着，快快地编译了休息；倘肯多花一文，便可以买一碟调试用例，或者跑一段开心代码，做练手项目了；如果能出到十几文，那就可以买一个完整项目。但这些顾客，多是短衣帮，大抵没有这样阔绰。只有穿格子衫的，才踱进店面隔壁的房子里，要项目要算力，慢慢地坐写。

我从十二岁起，便在镇口的咸亨码栈里当伙计。掌柜说，样子太傻，怕侍候不了格子衫主顾，就在外面做点事罢。外面的短衣主顾，虽然容易说话，但唠唠叨叨缠夹不清的也很不少。他们往往要亲眼看着代码从沙箱里拉出，看过你的代码有没有水，又亲看将沙箱放到服务器上，然后放心。在这严重监督下，掺水也很为难。所以过了几天，掌柜又说我干不了这事。幸亏荐头的情面大，辞退不得，便改为专管跑代码的一种无聊职务了。

我从此便整天地站在柜台里，专管我的职务。虽然没有什么失职，但总觉得有些单调，有些无聊。掌柜是一副凶脸孔，主顾也没有好声气，教人活泼不得；只有孔乙己到店，才可以笑几声，所以至今还记得。

孔乙己是站着接单而穿格子衫的唯一的人。他身材很高大；青白脸色，皱纹间时常夹些伤痕；一部乱蓬蓬的花白胡子。穿的虽然是格子衫，可是又脏又破，似乎十多年没有补，也没有洗。他对人说话，总是满口捕获借用，教人半懂不懂的。因为他姓孔，别人便从描红纸上的“上大人孔乙己”这半懂不懂的话里，替他取下一个绰号，叫作孔乙己。

孔乙己一到店，所有写代码的人便都看着他笑。有的叫道：“孔乙己，你头顶又空出一块了！”

他不回答，只对柜里说：“开两个沙箱，跑一段开心代码。”说完，排出九文大钱。

他们又故意高声嚷道：“你一定又写返回闭包了！”

孔乙己睁大眼睛说：“你怎么这样凭空污人清白……”

“什么清白？我前天亲眼见你写了个返回闭包，被编译器报 MGL1204，吊着打。”

孔乙己便涨红了脸，额上的青筋条条绽出，争辩道：“借用不能算逃逸……闭包！……程序员的事，能算逃逸么？”

接着，便是一串难懂的话，什么“const\& 只读”，什么“mut\& 写回”之类，引得众人都哄笑起来。店内外充满了快活的空气。

听人家背地里谈论，孔乙己原来也读过书，还钻研过语言设计，但终于没有当程序员，又不会营生。于是愈过愈穷，弄到将要断网了。幸而写得一手好代码，便替人家敲敲代码，接点外包，换一点托管费。

可惜他又有一样坏脾气，便是好闲懒做。坐不到几天，便连人和电脑键盘鼠标，一齐失踪。如是几次，叫他接外包的人也没有了。

孔乙己没有法，便免不了偶然做些偷窃的事。但他在我们店里，品行却比别人都好，就是从不拖欠。虽然间或没有现钱，暂时记在粉板上，但不出一月，定然还清，从粉板上拭去了孔乙己的名字。

孔乙己跑过半段代码，涨红的脸色渐渐复原。旁人便又问道：“孔乙己，你当真会闭包么？”

孔乙己看着问他的人，显出不屑置辩的神气。

他们便接着说道：“你怎的连半个 counter 也捞不到呢？”

孔乙己立刻显出颓唐不安模样，脸上笼上了一层灰色。嘴里说些话，这回全是“默认捕获是只读借用”之类，一句也听不懂了。

这时候，众人也都哄笑起来。店内外充满了快活的空气。

在这些时候，我可以附和着笑，掌柜是决不责备的。而且掌柜见了孔乙己，也每每这样问他，引人发笑。

孔乙己自己知道不能和他们谈天，便只好向孩子说话。

有一回，他对我说道：“你读过书么？”

我略略点一点头。

他说：“读过书……我便考你一考。闭包的捕获方式，有几种？”

我想，连沙箱都开不起的人，也配考我么？便回过脸去，不再理会。

孔乙己等了许久，很恳切地说道：“不能写罢？……我教给你，记着！这些方式应该记着。将来写回调的时候，要用。”

我暗想，我和写回调的还差得远呢。再说，我们掌柜也从不把闭包上账。我又好笑，又不耐烦，懒懒地答他：“谁要你教，不是四种么？”

孔乙己显出极高兴的样子，将两个指头的长指甲敲着柜台，点头说：“对呀对呀！……闭包有四种捕获，你知道么？”

我愈不耐烦了，努着嘴走远。

孔乙己刚用指甲在柜台上画 fn mut()，见我毫不热心，便又叹一口气，显出极惋惜的样子。

有几回，邻舍孩子听得笑声，也赶热闹，围住了孔乙己。他便给他们一人一段Magel代码。

孩子跑完示例，仍然不散，眼睛都望着屏幕。

孔乙己着了慌，伸开五指将屏幕罩住，弯腰下去说道：“不多了，我已经不多了。”

直起身又看一看屏幕，自己摇头说：“不多不多！多乎哉？不多也。”

于是这一群孩子都在笑声里走散了。

孔乙己是这样的使人快活，可是没有他，别人也便这么过。

有一天，大约是中秋前的两三天，掌柜正在慢慢结账，取下粉板，忽然说：“孔乙己长久没有来了。还欠十九个钱呢——那还是违约金零头。”

我才也觉得他的确长久没有来了。

一个写代码的人说道：“他怎么会来？……他写代码写垮了。”

掌柜说：“哦！”

“他总仍旧是接外包。这一回，是自己发昏，竟接了个写编译器的大项目。那家的活，是人接的么？”

“后来怎么样？”

“怎么样？先写文档，后来是改，改了大半夜，越改越写不出来。”

“那家的活，难在哪？”

“难在哪？闭包捕获要全过，借用检查要全过，循环引用检测要全过，还要零警告。少一样，都不算交付。”

“后来呢？”

“后来交不出货了。”

“交不出货怎样呢？”

“怎样？……合同上写着违约金，许是赔不起了。”

掌柜也不再问，仍然慢慢算他的账。

中秋过后，秋风一天凉比一天。看看将近初冬，我整天靠着火，也须穿上棉袄了。

一天的下半天，没有一个顾客，我正合了眼坐着。忽然间听得一个声音：“开一个沙箱。”

这声音虽然极低，却很耳熟。看时又全没有人。

站起来向外一望，孔乙己便在柜台下，对着门槛坐着。

他脸上黑而且瘦，已经不成样子；穿一件破夹袄，怀里抱着个旧键盘，手指还搭在键帽上；见了我，又说道：“开一个沙箱。”

掌柜也伸出头去，一面说：“孔乙己么？你还欠十九个钱呢——违约金零头还没清呢！”

孔乙己很颓唐地仰面答道：“这……下回还清罢。这一回是现钱，代码要新。”

掌柜仍然同平常一样，笑着对他说：“孔乙己，你又写循环引用了！”

但他这回却不十分分辩，单说了一句：“不要取笑！”

“取笑？要是不写，怎么连回收器都不拦你，跑出来个 counter 只涨不跌？”

孔乙己低声说道：“循环引用，对象里存闭包，闭包又捕获对象，环套环，回收不掉的……”

他的眼色，很像恳求掌柜，不要再提。

此时已经聚集了几个人，便和掌柜都笑了。

我开了沙箱，放到门槛上。

他从破衣袋里摸出四文大钱，放在我手里。见他满手是茧，指节上还有键帽印，原来，他便用这手敲了一路的代码。

不一会，他跑完代码，便又在旁人的说笑声中，抱着旧键盘，慢慢走去了。

自此以后，又长久没有看见孔乙己。到了年关，掌柜取下粉板说：“孔乙己还欠十九个钱呢——违约金零头。”

到第二年的端午，又说：“孔乙己还欠十九个钱呢——违约金零头。”

到中秋可是没有说，再到年关也没有看见他。

我到现在终于没有见——大约孔乙己的确死了。

\subsection*{技术对照}

返回闭包本身合法；问题是它是否携带活得不够久的借用。默认只读借用、mut 写回借用、copy 拥有副本和 move 转移拥有关系各有规则。循环持有与借用逃逸也不是同一个问题，应按实际所有权图和宿主资源协议分析，不能从人物对白推断存在通用追踪式回收器。
